% !TEX program = xelatex
\documentclass[UTF8,a4paper,zihao=5,fontset=windows]{ctexart}
\usepackage[margin=2.7cm]{geometry}
\usepackage{booktabs,enumitem,fancyhdr,array,needspace}
\usepackage{xurl}
\usepackage[hidelinks,unicode]{hyperref}
\setlength{\parindent}{2em}
\setlength{\parskip}{.28em}
\setlength{\emergencystretch}{3em}
\linespread{1.12}
\setlist{nosep,leftmargin=2em}
\renewcommand{\arraystretch}{1.16}
\newcolumntype{L}[1]{>{\raggedright\arraybackslash}p{#1}}
\pagestyle{fancy}\fancyhf{}\fancyfoot[C]{\thepage}\renewcommand{\headrulewidth}{0pt}
\hypersetup{pdftitle={AI工具使用详情},pdfauthor={},pdfsubject={全国大学生数学建模竞赛支撑材料}}
\newcommand{\tabnote}[1]{\par\vspace{0.3em}{\small #1\par}}
\begin{document}
\begin{center}{\LARGE\bfseries AI工具使用详情}\par\medskip
{\large《考虑变物性与径向收缩的圆柱药材干燥模型》}
\end{center}

本参赛队在竞赛过程中使用了AI工具。核心模型建立与论文初稿由参赛队完成；AI仅辅助求解代码编写、文献检索和语言润色。AI参与完成的代码与文献均经队员人工审查与核实后采用。

\section{所用AI工具名称、版本或型号}
\noindent
\begin{tabular}{@{}L{.26\linewidth}L{.18\linewidth}L{.49\linewidth}@{}}
\toprule
工具名称及提供方 & 版本或型号 & 使用方式 \\\midrule
ChatGPT桌面端（OpenAI）&
GPT-6 Astra&
竞赛期间使用该桌面端对话。\newline 未使用其他生成式AI平台，也未调用独立代码补全插件。\\
\bottomrule
\end{tabular}

\tabnote{Python、NumPy、SciPy、Matplotlib、openpyxl及XeLaTeX属于本地执行与排版工具链，单独运行这些程序不属于生成式AI。AI辅助编写的是调用上述工具链的求解、导出与核验脚本；实际计算在队员本机完成。}

\section{具体使用目的和环节}
核心建模与分析由参赛队主导。各环节分工如下。

\noindent
\begin{tabular}{@{}L{.16\linewidth}L{.42\linewidth}L{.35\linewidth}@{}}
\toprule
环节 & 参赛队完成的工作 & AI辅助的具体目的 \\\midrule
核心建模 & 确定四问有效传热传质框架、假设、控制方程、边界条件、环境延拓、收缩运动学及达标判据。 & 不用于生成或替代核心模型。\\
论文初稿 & 完成问题重述、问题分析、模型叙述、结果解释、结论及初稿结构。 & 不用于撰写初稿，仅在成稿后作语言润色。\\
问题一代码 & 给定常物性、固定半径与30分钟输出要求，审定离散方案。 & 辅助编写径向求解、表1和表2导出及数值检查脚本。\\
问题二代码 & 给定变物性局部耦合、前3小时规定表及全过程输出要求。 & 辅助编写变系数联立求解、1秒网格导出及与前3小时工作簿的核对。\\
问题三代码 & 给定4小时后环境延拓，以及全药材最大干基含水率严格低于0.15的判据。 & 辅助编写事件检测、极值复核、严格执行点选取及表5导出。\\
问题四代码 & 给定附件2半径历程、材料坐标与当前区域输出要求。 & 辅助编写收缩域求解、域外空白处理、表6导出及几何/质量检查。\\
文献检索 & 确定需要核对的方法类型，阅读原文并决定引用范围。 & 辅助检索圆柱热质耦合、动边界与数值积分等方面的文献线索。\\
语言润色 & 确认润色不得改动模型、公式、数字和结论。 & 调整摘要、章节衔接与个别句子的表达。\\
\bottomrule
\end{tabular}

\tabnote{图表由已核验求解结果的后处理脚本生成，属于求解代码的配套输出，不作为独立的AI绘图环节。定量图来自保存的数值数组或题表，流程图只表达队员已确定的求解步骤。}

\section{主要提示方式与使用过程说明}
主要提示方式为：先给出题面条件、队员已确定的模型方程和输入数据，再明确本次只要代码、文献线索或文字润色，并列出不得改动的边界。随后根据编译或运行报错、数值检查结果或原文核对意见，进行多轮交互修改，直到队员确认可以采纳。

以下为整理后的\textbf{典型交互示例}，用于说明提示方式和过程。

\subsection{求解代码：严格达标时长}
\textbf{输入材料：}问题二已采纳的变物性求解轨迹；附件1在0--4小时的环境观测；队员确定的4小时后末段平台延拓；判据为全药材最大干基含水率严格小于0.15。

\textbf{提示要点：}“按已确定的一维径向模型继续积分。不要用表格四位舍入值判断是否达标。先识别下降穿越0.15的等号根，再给出满足严格不等式的执行时刻，并输出该时刻的径向分布。”

\textbf{交互过程：}首轮代码能够积分并写出工作簿，但存在用舍入值判定、把等号根直接当作执行点、或未复核径向最湿位置等问题。队员指出后，要求把阈值根、极值搜索与执行点分开，并补网格/积分对照。修改后再本地运行，核对未舍入最大值确实小于0.15。

\textbf{采纳结果：}问题三采用一维条件执行时长57.4741小时。问题四按同样方式处理收缩域，采用51.0921小时；固定空间列落在当前半径以外时保留空白。

\subsection{求解代码：前三问离散与收缩实现}
问题一、二的提示同样先给出队员已定方程，再要求编写可运行求解器。典型要求包括：温度与干基含水率联立；中心对称、表面第三类边界；物性按附录在单元内更新；输出规定时刻和规定半径。问题四另要求在材料坐标下推导，区分材料速度与网格速度，不能只把常数半径替换为$R(t)$。

AI给出的代码经队员对照方程修改后，在本机用SciPy积分器运行。问题一采用高分辨率有限体积；问题二、三采用统一变物性轨迹；问题四采用材料坐标收缩求解。规定表1--6及第二问全过程1秒网格由上述程序导出。

\subsection{文献检索}
\textbf{输入材料：}圆柱药材干燥、热质耦合、第三类边界、有限体积或动边界等关键词，以及本题已采用的有效物性和收缩数据范围。

\textbf{提示要点：}“列出可核对原文的圆柱或近圆柱干燥热质传递文献，说明其离散、边界或运动学与本题可能对应的部分，不要根据摘要直接抄公式。”

\textbf{交互过程：}AI返回若干候选。队员随后阅读原文，核对比热、扩散、边界单位和收缩定义是否与本题变量一致，只引用实际核对过、且未超出原文适用范围的内容。

\textbf{采纳结果：}正文参考文献中的da Silva等（2014）、Adrover等（2020）、吴孟秋等（2022）用于圆柱控制体、材料运动与动网格组织参考；SciPy手册用于积分器说明；ASHRAE与FAO文献用于环境量换算参考。未把各文的香蕉、梨或白萝卜本构参数写入本题主模型。

\Needspace{9\baselineskip}
\subsection{语言润色}
\textbf{输入材料：}队员已完成的摘要、章节过渡句或图注草稿，并附带“不得改动任何数字、符号、表号和结论指向”的限制。

\textbf{提示要点：}只调整语气、衔接和重复表述，保持原有加粗范围与结果含义。

\textbf{交互过程：}对返回文字逐句对照原稿。若出现改变时长含义、把条件结果写成实测保证、或改动表格数字的情况，一律改回。语言润色本身不产生新的模型或数值。

\section{对AI输出的采纳、人工修改和核验}
本队对求解代码与文献检索均作了人工审查；语言润色只核对未改动实质内容。

\noindent
\begin{tabular}{@{}L{.16\linewidth}L{.42\linewidth}L{.35\linewidth}@{}}
\toprule
内容 & 采纳与人工修改 & 核验情况 \\\midrule
求解代码 &
采纳经修改后的四问求解、导出和检查脚本。主要修改：边界通量与队员方程一致，输出时域与题表对齐，达标判定改用未舍入量，收缩域外留空，并拒绝把含潜热试验代码混入主程序。&
队员阅读离散格式并对照已定方程；在本机实际运行；核对表1--6与对应工作簿；问题三、四检查未舍入$\max C<0.15$；并用独立网格、积分方法或有限体积对照限制结论范围。\\
文献检索 &
采纳已阅读原文、且与本题方法对应的6篇参考文献。不采纳仅有摘要、单位不符或收缩定义不同的条目，也不把原文经验系数写入主模型。&
队员核对外文公式符号、单位和适用对象；正文引用与参考文献一一对应。\\
数值结果 &
采用问题一现行30分钟结果，问题二补齐后的全过程轨迹，问题三57.4741小时和问题四51.0921小时的一维执行点。&
结果来自队员运行后的程序输出，不以对话中的口述数字为准。二维与后段环境只作为已有检查或条件情景，不改写上述正式时长。\\
语言润色 &
仅采纳未改变数字、公式和结论的表述调整。&
队员通读确认含义未变。\\
\bottomrule
\end{tabular}

\vspace{0.6em}
未采纳或已纠正的代表性输出包括：把表格舍入值当作严格达标依据；把等号根直接当作执行时刻；只把常数半径换成$R(t)$而不改方程；把文献中的潜热项、等温线或经验收缩率直接写入本题；把条件对比写成实际节能率。对AI输出均以队员已定模型和实际运行结果为准，不因回复完整而直接采用。


\section{支撑包复现链补充记录}
本轮根据“支撑材料中的代码能否复现论文结果”及“你来完成”的请求，补充独立工作目录中的实际数值复跑、原始输入与源码哈希核对、常规Excel导出、逐格比较以及运行说明。保留原模型和数值算子，新增复现控制与核对代码。

原题设表、临界根和完整提交数据分别按显示精度与记录的数值容差核对。程序读回导出工作簿，并以故意改错的样本检查比较器能够报警。实际结果、运行命令与退出状态存于复现验收记录。

以上为工具实际执行的技术核查；新增代码及记录的最终采纳和人工审阅由参赛队负责，自动检查不替代人工确认。原文模型适用范围和第四问二维全域验证限制保留。

\end{document}
